\documentclass[11pt]{article}
\usepackage{amsmath, amssymb}
\usepackage{geometry}
\usepackage{enumitem}
\usepackage{hyperref}
\geometry{margin=1in}

\title{\textbf{CSE400 – Fundamentals of Probability in Computing}}
\author{\textbf{Moksha Sheth} \\\textbf{Lecture Scribe-2}}

\date{February 11 $2026$}

\begin{document}
\maketitle

\section*{Lecture 10: Randomized Min-Cut Algorithm (February 5, 2026)}\\

\hrule
\vspace{1.5cm}

\section{Overview of the Lecture}
This lecture focuses on the Min-Cut Problem in graphs and builds toward Karger’s Randomized Min-Cut Algorithm, while also covering a deterministic alternative (Stoer--Wagner). The flow is:
\begin{itemize}
\item Min-Cut Problem
\item Why use min-cut?
\item What is min-cut? (including edge contraction)
\item Successful vs unsuccessful min-cut runs
\item Max-Flow Min-Cut Theorem
\item Deterministic min-cut (Stoer--Wagner) + pseudocode
\item Randomized min-cut (Karger) + pseudocode
\item Comparison: deterministic vs randomized
\item Theorem about probability of outputting a min-cut set
\item Python simulation note
\end{itemize}

\vspace{1cm}
\section{Key Definitions and Terminology}

\subsection*{\underline{Cut-set (in a graph)}}
A cut-set is a set of edges whose removal breaks the graph into two or more connected components.

\subsection*{\underline{Minimum cut (Min-cut) problem}}
Given a graph $G = (V, E)$ with $n$ vertices, the minimum cut problem is to find a minimum cardinality cut-set in $G$.

\subsection*{\underline{Edge contraction}}
The main operation used (especially in randomized min-cut): edge contraction removes an edge while simultaneously merging its two endpoints.\\
Contracting an edge $(u, v)$:
\begin{itemize}
\item merge $u$ and $v$ into one vertex,
\item eliminate all edges connecting $u$ and $v$,
\item retain all other edges.
\end{itemize}
The new graph may have parallel edges but no self-loops.

\subsection*{\underline{Successful min-cut run}}
A run where the algorithm successfully finds the minimum cut of the graph.

\subsection*{\underline{Unsuccessful min-cut run}}
A run where the algorithm fails to correctly identify the minimum cut.

\subsection*{\underline{Flow-network terms (from Max-Flow Min-Cut)}}
Capacity of a cut: sum of capacities of edges in the cut oriented from a vertex in $X$ to a vertex in $Y$.\\
Minimum cut: cut with the smallest possible capacity.\\
Minimum cut capacity: capacity of the minimum cut.\\
Maximum flow: largest possible flow from source $S$ to sink $T$.

\vspace{1cm}
\section{Explanation of Core Concepts and its Understanding}

\subsection*{A. \underline{Why use Min-Cut? (Applications mentioned)}}
Min-cut is used in applications related to network connectivity, reliability, and optimization, including:
\begin{itemize}
\item Network Design: find the minimum capacity cut to improve communication efficiency and optimize flow.
\item Communication Networks: analyze vulnerability to failures; helps build robust, fault-tolerant networks.
\item VLSI Design: partition circuits into smaller components to reduce interconnectivity complexity.
\end{itemize}
(Reference mentioned in slides: Section 1.5 ``Application: A Randomized Min-Cut Algorithm'' in \textit{Probability and Computing} (2nd Edition).)

\subsection*{B. \underline{What is Min-Cut and why randomness matters?}}
Min-cut algorithms like Karger’s are random and can be sensitive to early choices. If the algorithm contracts ``critical edges'' early, it can affect whether it finds the true minimum cut.

\subsection*{C. \underline{Edge contraction as the core operation}}
Edge contraction repeatedly reduces the graph until only two super-nodes remain; the edges between them represent a cut. Contraction can introduce parallel edges, and self-loops are removed.

\subsection*{D. \underline{Successful vs Unsuccessful runs (idea)}}
The lecture visually illustrates that some contraction choices preserve the true min-cut until the end (successful), while other random contraction sequences destroy it early (unsuccessful).

\subsection*{E. \underline{Max-Flow Min-Cut Theorem (core theorem statement)}}
The max-flow min-cut theorem: In a flow network, the maximum amount of flow from source to sink equals the total weight of edges in a minimum cut.

\vspace{1cm}
\section{Important Formulas and Mathematical Expressions}

\subsection*{\underline{Max-Flow Min-Cut Theorem (statement form used in lecture)}}
``In a flow network, the maximum amount of flow passing from the source to the sink is equal to the total weight of the edges in a minimum cut''.

\subsection*{\underline{Time complexity statements (as given)}}
Stoer--Wagner deterministic min-cut:
\[
O\!\left(V \cdot E + V^{2}\log V\right)
\]
where $V$ = number of vertices, $E$ = number of edges.\\[4pt]
Karger’s randomized min-cut:
\[
O\!\left(V^{2}\right)
\]
where $V$ = number of vertices.

\subsection*{\underline{Probability theorem for min-cut set (as shown)}}
The algorithm outputs a min-cut set with probability at least:
\[
\frac{2}{n(n-1)}
\]
(with $n$ vertices).

\vspace{1cm}
\section{Derivations / Proofs (If Present in Lecture)}
No step-by-step derivations were presented in the lecture slides; the lecture provides theorem statements and algorithmic reasoning (e.g., Max-Flow Min-Cut theorem statement, Stoer--Wagner statement, and the probability guarantee statement).

\vspace{1cm}
\section{Worked Examples / Applications}

\subsection*{A. \underline{Application contexts (from ``Why use min-cut?'')}}
\begin{itemize}
\item Network Design
\item Communication Networks
\item VLSI Design
\end{itemize}

\subsection*{B. \underline{Example-style illustrations shown in lecture}}
Edge contraction illustration: shows how contracting an edge merges vertices and updates the graph (parallel edges possible; no self-loops).\\
Successful min-cut run figure: shows a contraction sequence that ends with the correct min-cut.\\
Unsuccessful min-cut run figure: shows a contraction sequence that fails to preserve the minimum cut.\\
(These are presented as conceptual/visual examples in the slides rather than numeric step problems.)

\vspace{1cm}
\section{Important Observations and Results}

\subsection*{\underline{Deterministic Min-Cut (Stoer--Wagner) -- Key result stated}}
Let $s$ and $t$ be two vertices of a graph $G$. Let $G/\{s,t\}$ be the graph obtained by merging $s$ and $t$. Then a minimum cut of $G$ can be obtained by taking the smaller of:
\begin{itemize}
\item a minimum $s$-$t$ cut of $G$, and
\item a minimum cut of $G/\{s,t\}$.
\end{itemize}
Reason given (as stated):\\
Either there is a minimum cut separating $s$ and $t$ $\rightarrow$ then min $s$-$t$ cut is min cut of $G$.\\
Or there is none $\rightarrow$ then min cut of $G/\{s,t\}$ works.

\subsection*{\underline{Why Randomized Algorithms? (points listed)}}
Randomized algorithms provide a probabilistic guarantee of success and can give a more accurate estimate with fewer iterations; listed motivations include:
\begin{itemize}
\item Efficiency
\item Parallelization
\item Approximation guarantees
\item Avoidance of worst-case instances
\item Heuristic nature
\item Robustness
\end{itemize}

\subsection*{\underline{Comparison summary (as given)}}
Deterministic min cut: exact guarantee; may be higher time complexity for large graphs.\\
Randomized min cut: approximate min cut with high probability; lower time complexity statement given for Karger.\\
``Any specific problem is responsible for deciding the approach.''

\subsection*{\underline{Python simulation note (lecture instruction)}}
Open Campuswire post for Lecture 10 and download the provided \texttt{.ipynb} file from comments.

\vspace{1cm}
\section{Quick Revision Points}
\begin{itemize}
\item Min-cut: minimum cardinality cut-set of $G = (V, E)$.
\item Cut-set: edges whose removal disconnects the graph.
\item Core operation (randomized min-cut): edge contraction (merge endpoints; remove self-loops; parallel edges may occur).
\item Successful vs unsuccessful runs: contraction choices can preserve or destroy the true min-cut.
\item Max-Flow Min-Cut theorem: max flow equals capacity (total weight) of a min cut.
\item Stoer--Wagner: deterministic approach; uses repeated phases and merging based on the minimum cut-of-the-phase.
\item Karger’s randomized algorithm: uses random contractions; provides probabilistic success guarantee.
\item Complexities: Stoer--Wagner $O(V \cdot E + V^{2}\log V)$, Karger $O(V^{2})$.
\item Probability guarantee: outputs a min-cut set with probability at least $\dfrac{2}{n(n-1)}$.
\end{itemize}

\end{document}